\documentclass[../main.tex]{subfiles}
\begin{document}
%==========================================TIÊU ĐỀ TẬP========================================
\begin{table}[h]
  \begin{adjustwidth}{-1cm}{}
    \begin{tabular}{>{\centering\arraybackslash}p{6cm}|>{\raggedright\arraybackslash}p{12.5cm}}
      \multirow{5}{6cm}
      % Cột bên trái: ÉPISODE và số tập
      {\\[-40pt]
        \flaregothic\fontsize{20pt}{20pt}\selectfont \centering
        \color{eptype}ÉPISODE\color{black}\\
        \burbank\fontsize{72pt}{72pt}\selectfont
        \color{epnum}24\color{black}\\[6pt]
        \cafeteria\fontsize{20pt}{20pt}\selectfont\color{epnum}(2-7)\color{black}
        \\[18pt]
        \flaregothic\fontsize{14pt}{14pt}\selectfont
        \color{epattr}BIENVENUE, \uline{\color{Okanya} Okayu} !
        \flaregothic\fontsize{18pt}{18pt}\selectfont
        \color{epattr}ÉPISODE PERDU
        \color{black}
      }
       &
      % Dòng thứ nhất: tên tập tiếng Nhật
      {\fotskip\fontsize{18pt}{18pt}\selectfont \ruby{変}{へん}なヤツがいるんだけどwwwww
      }          \\[6pt]
      % Dòng thứ hai: tên tập tiếng Anh
       & {\cafeteria\fontsize{18pt}{18pt}\selectfont OkaKoro and the Weird Visitor} \\[4pt]
      % Dòng thứ ba: tên tập tiếng Việt
       & {\vnmsans\fontsize{18pt}{18pt}\selectfont Bóng đen bí ẩn}                  \\[8pt]
      % Dòng thứ tư: Nhân vật xuất hiện
       & {\caviar{\fontsize{14pt}{14pt}\selectfont Nhân vật: \emp{kuon1} \emp{ryuuhou1} \emp{achan} \emp{mel} \emp{okayu} \emp{korone}}} \\[8pt]
      % Dòng cuối cùng: Thời gian diễn ra của tập (thường sẽ là ngày phát hành)
       & {\continuum\fontsize{16pt}{16pt}\selectfont Ngày phát hành: 20/10/2019}    \\[18pt]
    \end{tabular}
  \end{adjustwidth}
\end{table}
%=========================================TRUYỆN CHÍNH========================================
\color{black}

\txtbx[2]{{\color{vol} \uline{Tóm tắt:}} Một bóng đen bí ẩn xuất hiện trong văn phòng và bắt chước mọi hành động của Korone, giúp cô nàng thức tỉnh năng lượng Stand đầy quyền năng. Thế nhưng, sức mạnh mới này lại vô tình đẩy Okayu vào một chuyến hành trình ngoài không gian bất đắc dĩ, nơi cô nàng phải đối mặt với thảm họa thiên thạch để giải cứu Trái Đất.}

\model{24}{r}{6cm}{okayu}

Trong không gian yên tĩnh của văn phòng, nàng Khuyển Korone đang ngồi bệt giữa sàn nhà, ánh mắt dán chặt vào một góc tường với vẻ đầy cảnh giác. Dù bầu không khí xung quanh chẳng có gì bất thường, nhưng bản năng loài chó mách bảo cô rằng có một sự hiện diện vô cùng kỳ quái đang lẩn khuất đâu đó. 

Bất thình lình, một giọng nói nhẹ nhàng, mang theo chút âm hưởng lười biếng nhưng vô cùng quen thuộc vang lên ngay phía sau: ``Koro-san\~{} Đang làm gì mà chăm chú thế?''

\chrint

Đó là Nekomata Okayu (\ruby{猫}{ねこ}\ruby{又}{また}おかゆ, \ruby{Nê-cô}{Miêu}-\ruby{ma-ta}{Hữu}\ Ô-ca-iu), cô tiểu thư của một gia đình làm cơm nắm danh tiếng tại Tokyo. Là thành viên cuối cùng của nhóm \hll\ GAMERS, Okayu nổi tiếng với mối quan hệ ``trên mức tình bạn" với Korone, tạo thành cặp bài trùng Oka-Koro (hay Chó-Mèo) lừng danh khắp công ty.

Okayu sở hữu tính cách phóng khoáng, tự do đến mức đôi khi người ta thấy cô nàng thật ``phởn". Tuy nhiên, đằng sau vẻ ngoài thư thái đó lại là một kẻ thích chơi khăm và có thói quen ``tán tỉnh" dạo với các thành viên khác chỉ để xem phản ứng của họ. Cô nàng cũng được biết đến là một trong những người có đầu óc ``đen tối" và lém lỉnh nhất hololive, luôn biết cách khuấy động bầu không khí bằng những chiêu trò của mình.

Dù vậy, Okayu cũng là một tâm hồn nhạy cảm. Cô có thể bật khóc nức nở khi chứng kiến những khoảnh khắc xúc động trong các tựa game mình chơi và cần một khoảng thời gian yên lặng để lấy lại bình tĩnh. Với mái tóc tím nhạt, đôi tai mèo tinh nghịch và chiếc áo hoodie đen in hình cơm nắm đặc trưng, Okayu luôn toát lên một sức hút kỳ lạ, vừa gần gũi vừa đầy bí ẩn.

\chrint

\img{2-8a}

Okayu bước đến gần, miệng nở một nụ cười tươi tắn, ánh mắt cũng tò mò dõi theo hướng mà Korone đang nhìn. Cô hỏi khẽ: ``Bà đang rình con gì à?''

Nàng Khuyển khẽ ngoảnh lại nhìn người bạn thân nhất của mình, rồi run run chỉ tay về phía góc tối, giọng đầy vẻ dè chừng: ``Có cái gì đó\dots\ kỳ lạ lắm đang ở đằng kia kìa Okayu.''

\img{2-8b}

Nơi Korone chỉ là một vật thể đen kịt đang không ngừng chuyển động một cách chậm chạp. Nó không giống như bóng của đồ vật bình thường, mà dường như có ý thức riêng, liên tục thay đổi hình dạng và lướt đi trong không gian mờ ảo của căn phòng.

Ngồi gần đó, Ryuuhou đang cầm máy game cầm tay cũng phải ngẩng lên. Đôi mắt hai màu của cô bé mở to, đầy vẻ phấn khích thay vì sợ hãi: ``Onii-san, cái bóng đó là do anh mang về à?''

Kuon đang vùi đầu vào đống báo cáo thu chi, nghe tiếng em gái gọi mới mệt mỏi ngước lên. Nhìn thấy cái bóng đen đang ngoe nguẩy, cậu không giấu nổi sự e dè: ``Anh cũng chẳng biết nó là cái giống gì nữa. Nó cứ lù lù xuất hiện ở đó ngay từ khi Korone-chan bước vào văn phòng rồi.''

Ba người cùng tiến lại gần để quan sát kỹ hơn thực thể kỳ quái này. Okayu thốt lên đầy thích thú: ``Kinh thật! Nhìn nó cứ như cục thạch đen vậy!''

Korone thì nhận ra một điều thú vị khác: ``Này, tôi phát hiện ra là cứ mỗi khi tôi làm động tác gì, nó lại bắt chước theo y hệt luôn này!''

\img{2-8c}

Để chứng minh, nàng Khuyển dang rộng hai tay ra. Quả nhiên, bóng đen kia cũng vươn ra hai phần hình thù giống hệt đôi tay, dù ánh sáng yếu ớt khiến mọi thứ trông thật mờ ảo. 

Ryuuhou thích thú vỗ tay: ``Hay quá! Giống như Stand trong JoJo ấy! Onii-san nhìn kìa, nó còn vẫy tay lại với chị ấy nữa này!''

Kuon nhíu mày, thử vẫy tay về phía cái bóng nhưng đáp lại cậu chỉ là sự im lặng tuyệt đối. Cái bóng đen hoàn toàn không có chút phản ứng nào với cậu quản lý. Cậu thở dài: ``Chán thật, sao nó chẳng thèm để ý gì đến anh thế này? Có vẻ như nó chỉ trung thành với mỗi Koro-chan thôi thì phải.''

Okayu cười khúc khích: ``Chắc tại Kuon-senpai không có `tố chất' làm người huấn luyện chó rồi!''

Thế là cả căn phòng rộn rã tiếng cười khi Korone bắt đầu thử đủ mọi trò nghịch ngợm với cái bóng, từ nhảy múa cho đến trò trượt trên sàn nhà với những con cá thu đông lạnh thật sự. 

\img{2-8d}

Tuy nhiên, sau một hồi vui đùa, sự tò mò lại trỗi dậy. Kuon trầm ngâm hỏi: ``Nhưng rốt cuộc thì cái thứ này từ đâu mà ra nhỉ? Chẳng lẽ có ai đó đã gửi nó đến đây sao?''

\img{2-8e}

Korone hướng mắt ra phía cửa sổ, đưa ra một giả thuyết đầy màu sắc khoa học viễn tưởng: ``Có khi nó là sinh vật ngoài hành tinh, đi theo một mảnh thiên thạch đâm xuống Trái Đất của chúng mình thì sao?''

Cả nhóm cùng quay ra nhìn qua cửa sổ. Cảnh tượng bên ngoài bỗng chốc trở nên vô thực như trong một bộ phim giả tưởng: bầu trời rực lên sắc vàng cam kỳ dị, và một mặt trăng khổng lồ đến mức bất thường đang tỏa sáng rực rỡ, soi rọi những ngọn đồi núi lạ lùng thay vì những tòa nhà quen thuộc của thành phố.

Kuon lẩm bẩm: ``Anh lại thấy nó giống như đến từ một chiều không gian khác hơn. Vấn đề là tại sao nó lại chỉ chọn Koro-chan để bắt chước chứ?''

Một sự im lặng bao trùm lấy không gian. Okayu mỉm cười, cố gắng làm dịu bầu không khí đang dần trở nên căng thẳng: ``Thôi nào, em nghĩ nó cũng hiền mà. Coi như là một con cừu lạc tự mò đến đây tìm bạn thôi.''

Kuon vẫn không khỏi lo lắng: ``Vẫn phải cảnh giác, ai biết nó có đột ngột biến thành quái vật ăn thịt người không chứ.''

Ryuuhou thì chẳng mảy may lo lắng, cô bé tinh nghịch dùng ngón tay cấu nhẹ vào cái bóng đen: ``Nếu nó là cừu lạc thì chúng ta phải tìm cách mua vui cho nó chứ, đúng không Okayu-san?''

Nàng Khuyển ngây ngô hỏi lại: ``Mua vui? Làm thế nào để nó vui được?''

Nàng Mèo đặt tay lên cằm suy tư: ``Hừm, thường thì điều gì sẽ khiến bà thấy vui nhất?''

Nàng Khuyển mắt suy nghĩ cực kỳ nghiêm túc, rồi thốt ra một câu khiến Kuon suýt ngã ngửa: ``\dots Chắc là được cho tăng ca không lương?''

Mặt cậu Tổng tái mét, mồ hôi hột chảy ròng ròng: ``Thôi, anh xin em! Nhìn gương chị A-san là đủ hiểu rồi, chẳng có gì vui vẻ ở cái việc đó đâu!''

Cả nhóm bất giác rùng mình khi tưởng tượng đến cảnh A-chan - cô nàng đeo kính tận tụy - đang phải cặm cụi bên màn hình máy tính giữa đêm khuya, đôi mắt thâm quầng và gương mặt đầy vẻ mệt mỏi bên đống hồ sơ chất cao như núi. 

\begin{figure}[H]
  \centering
  \begin{subfigure}[b]{0.45\textwidth}
    \centering
    \includegraphics[width=8cm]{../../../../Image/Book2/2-8f1.png}
  \end{subfigure}
  \quad
  \begin{subfigure}[b]{0.45\textwidth}
    \centering
    \includegraphics[width=8cm]{../../../../Image/Book2/2-8f2.png}
  \end{subfigure}
  \caption{``Tại sao đời tôi lại khổ thế này\dots?''}
\end{figure}

Tiếng cười của Okayu vang lên phá tan bầu không khí ảm đạm: ``Đúng là chỉ có Koro-san mới nghĩ ra được trò đó! Thôi, mua vui thì chắc là cùng nó chơi mấy trò vận động mạnh đi!''

Kuon thở phào nhưng vẫn không quên dặn dò: ``Làm gì thì làm, đừng có kéo anh vào mấy cái trò điên rồ của hai đứa là được.''

\ct

Sáng hôm sau.

Cánh cửa văn phòng bị đạp tung một tiếng \textbf{*RẦM!*} khiến Okayu đang ăn cơm nắm cũng phải giật mình, đánh rơi cả cục cơm. Korone bước vào với vẻ mặt đầy hăng hái: ``\textbf{Chào buổi sáng mọi người!}''

Kuon vẫn dán mắt vào máy tính, buông một câu lạnh lùng: ``Lại phá cửa, trừ lương em nha, Koro-chan.''

Nhưng rồi, Korone bỗng khựng lại, vẻ mặt trở nên hoảng hốt, nhưng không phải vì bị cô bị cậu Tổng phạt. Cô nàng bắt đầu nhìn quanh người mình, lúng túng kêu lên: ``Ơ\dots\ Ể!? Cái quái gì đang xảy ra với tôi thế này!?''

\img{2-8h}

Okayu đặt hộp cơm xuống, chạy lại gần rồi tròn mắt ngạc nhiên: ``Trời đất! Bà\dots\ bà vừa mới thức tỉnh sức mạnh Stand thật đấy à Koro-san?''

Kuon lúc này mới ngẩng đầu lên và chết lặng. Toàn thân Korone giờ đây đang hòa làm một với cái bóng đen tối qua, tỏa ra một luồng ánh sáng tím mờ ảo đầy quyền năng và có phần đáng sợ. 

Ryuuhou đứng bên cạnh hét lên đầy phấn khích: ``Korone-san đã thức tỉnh Stand rồi! Onii-san nhìn kìa, nó còn phát sáng nữa này!''

Kuon lắp bắp: ``Anh không biết gì về Stand cả, nhưng nhìn cái khí thế này thì đúng là em ấy đang hăng máu quá mức rồi\dots''

Korone hoảng loạn, vung tay múa chân cầu cứu: ``\textbf{Okayu! Cứu tôi với, tôi không điều khiển được nó!}''

\begin{figure}[H]
  \centering
  \begin{subfigure}[b]{0.45\textwidth}
    \centering
    \includegraphics[width=8cm]{../../../../Image/Book2/2-8i1.png}
  \end{subfigure}
  \quad
  \begin{subfigure}[b]{0.45\textwidth}
    \centering
    \includegraphics[width=8cm]{../../../../Image/Book2/2-8i2.png}
  \end{subfigure}
\end{figure}

Nhưng khi Korone vừa định chạm vào tay Okayu, một luồng sức mạnh khủng khiếp bộc phát. Cú chạm vô tình biến thành một cú đấm uy lực kinh người, làm vỡ tan tành cửa kính chống đạn và hất tung Okayu ra ngoài cửa sổ, bay thẳng lên trời như một quả tên lửa hành trình.

``\textbf{O-OKAYU!!!}'' Korone gào lên thảm thiết.

Kuon đứng đơ như tượng gỗ, nhìn cái lỗ hổng to đùng trên tường rồi nhìn sang Korone với ánh mắt kinh hoàng. Cậu lẩm bẩm bằng tiếng mẹ đẻ: ``{\color{red} V-Vãi cả chưởng\dots\ Mạnh thế này thì ai chơi lại!? Thôi, chuồn nhanh còn kịp!} Ryuuhou, chạy đi!!''

Không chần chừ một giây, Kuon lao thẳng ra cửa với vận tốc ánh sáng, mặc cho tiếng gọi của Korone vang vọng phía sau: ``Khoan đã! Đừng bỏ em lại một mình với cái đống này mà, Kuon-senpai!!''

Ryuuhou thì chẳng chạy, cô bé đứng lại, chăm chú nhìn vào chiếc tivi trong góc phòng đang phát bản tin khẩn cấp từ holoNEWS với sự xuất hiện của phóng viên Yozora Mel.

\img{2-8j}

``Nguy hiểm tột độ! Một thiên thạch khổng lồ đang lao thẳng về phía Trái Đất!'' nàng Ma cà rồng hốt hoảng hét vào máy quay, gương mặt áp sát màn hình đầy vẻ tuyệt vọng. ``Liệu có Messiah\footnote{Vị cứu tinh giải cứu một nhóm người, thường được biết đến nhiều trong kinh thánh Hebrew (Do Thái) trong tôn giáo Abraham. Ý của Mel đúng với nghĩa đen của nó, tức là ``vị cứu tinh''.} nào xuất hiện để giải cứu nhân loại chúng ta khỏi thảm họa diệt vong này không!?''

Korone đứng lặng người nhìn bản tin, tâm trí hoàn toàn trống rỗng giữa mớ hỗn độn: người bạn thân vừa bị mình đấm bay lên trời, và giờ là cả thế giới sắp bị hủy diệt. Cô chỉ biết thốt lên một tiếng ``Ồ'' hờ hững trong vô vọng.

\ct

``Mình vẫn đang xoay nè-nya\dots!!''

\gif{2-8k}{1}{42}

Giữa không gian vũ trụ bao la, Okayu đang quay tít mù sau cú đấm của Korone. Đột nhiên, cô nàng đập thẳng vào mảnh thiên thạch khổng lồ đang lao tới. 

``Ể!? Cái méo gì to đùng thế này!?'' Okayu sững sờ khi nhận ra mình đang cưỡi trên một hòn đá tử thần. Dù cố gắng bám trụ nhưng cô nàng vẫn liên tục bị trượt đi trong vô vọng. 

Giữa lúc ngàn cân treo sợi tóc, bóng đen bí ẩn của Korone bỗng dưng xuất hiện ngay cạnh Okayu, như thể nó vừa dịch chuyển tức thời qua hàng vạn dặm không gian. Nó bao phủ lấy toàn bộ thiên thạch và bắt đầu siết chặt lại.

\img{2-8l}

Dưới sức mạnh kinh hoàng của bóng đen, mảnh thiên thạch khổng lồ dần bị hấp thụ và thu nhỏ lại, cuối cùng chỉ còn là một viên đá nhỏ xíu phát sáng lấp lánh trên tay Okayu. Nàng Mèo tím ngơ ngác nhìn viên đá: ``Ể\dots\ vậy là mình vừa cứu cả thế giới thật đấy à?''

\ct

Tại văn phòng lúc này, A-chan vẫn đang ngồi thẫn thờ với con cá thu đông lạnh, đôi mắt lờ đờ như người mất hồn: ``Bữa tối cuối cùng của nhân loại là cá thu\dots''

Korone nhìn lên tivi, thấy Okayu đang rạng rỡ phát biểu từ không gian với viên đá trên tay. 

``Thật không thể tin nổi! Okayu-chan chính là Messiah của chúng ta! Thế giới đã bình yên trở lại rồi!'' Mel reo lên đầy phấn khích trên sóng truyền hình.

Korone nhảy cẫng lên vui sướng: ``Oa! Okayu vẫn còn sống! Chị ấy đỉnh quá!''

Kuon lúc này mới lén lút ló đầu vào văn phòng, thở phào nhẹ nhõm: ``Vãi cả siêu anh hùng\dots\ Hóa ra cái bóng đen đó cũng có ích phết chứ chẳng chơi.''

Chỉ có A-chan là thở dài chán nản, đặt con cá thu xuống bàn: ``Vậy là tôi vẫn phải tiếp tục tăng ca\dots\ Trời ơi, sao đời tôi khổ thế này\dots''

%==========================================TRUYỆN PHỤ=========================================
\omk

\begin{center}
  \uline{Hậu trường.}
\end{center}

``Và\dots\ cắt! OK, ổ-ổn rồi đấy!''

Công việc quay phim đã hoàn thành. Để quay được số \textit{holo no Graffiti} này, A-chan và Kuon đã huy động các thành viên chuyên ban nghệ thuật từ trường Xây dựng Kawachi đến giúp.

Orikawa Yamazu (\ruby{織}{おり}\ruby{川}{かわ}\ruby{山}{やま}\ruby{住}{ず}, \ruby{Ô-ri}{Chí}-\ruby{ka-oa}{Xuyên}\ \ruby{I-a-ma}{Sơn}-\ruby{dư}{Trụ}), bạn học cùng lớp của Kuon, nổi bật với khả năng tiếng Anh khá ổn nhưng bị tật nói lắp, hôm nay được chỉ định làm ``đạo diễn''. Dưới biệt danh ``nghiện ngập kelvin'' trong nhóm chat, cậu đang hết sức phấn khích khi được chỉ huy cả một đoàn quay phim, dù tật nói lắp của mình có ảnh hưởng đôi chút.

Ryuuhou cũng có mặt trong vai trò ``trợ lý hậu đài'' đắc lực, cô bé nhanh nhẹn chạy lại đưa nước cho Okayu, người vừa bước xuống từ cái lồng xoay với gương mặt tái mét vì chóng mặt.

``Chị Okayu giỏi quá! Xoay thế mà vẫn nói được thoại!'' cô bé mắt hai màu khen ngợi.

Okayu lảo đảo, may mà có Korone và anh chàng nhân viên Fukuyasu đỡ kịp: ``C-Cảm ơn em Ryuuhou-chan\dots\ n-nhưng mà chị sắp nôn tới nơi rồi\dots''

Ở một góc khác, Jotaro Hokuro (\ruby{承}{じょう}\ruby{太}{た}\ruby{郎}{ろ}\ruby{北}{ほく}\ruby{麓}{ろ}, \ruby{Giô}{Chửng}-\ruby{ta}{Thái}-\ruby{rô}{Lang}\ \ruby{Hô-cư}{Bắc}-\ruby{rô}{Lộc}), người chịu trách nhiệm viết kịch bản cho số này, vừa quan sát vừa buông lời bình. Biệt danh ``Người mà ai cũng biết là ai'' của cậu thực sự phù hợp với phong thái tự tin và kỳ quái.

Kuon nhận xét: ``Kịch bản ông viết lần này nhảm sịt thật đấy, nhưng mà theo kiểu hài hước đúng chất của công ty này.''

Hokuro gãi đầu cười tự tin: ``Tôi đã cố tình chắp vá thêm đoạn thiên thạch để tạo kịch tính mà lị!''

Trong khi đó, Korone đang lo lắng xoa bụng cho Okayu: ``Xin lỗi bà nha Ogayu, cú đấm nãy tôi có làm hơi quá tay không?''

Okayu cười khì khì, dù vẫn còn hơi choáng: ``Không sao đâu, hy sinh vì nghệ thuật mà! Miễn là fan thấy vui là được!''

Nhìn cặp đôi Chó-Mèo quấn quýt bên nhau, Ryuuhou huých tay anh trai mình: ``Nhìn họ thân thiết thật Onii-san, đúng là cặp bài trùng!''

Kuon mỉm cười, nhìn đống đổ nát (giả) trên phim trường: ``Ừ, dù có hơi loạn lạc nhưng mà cũng vui. Giờ thì dọn dẹp rồi đi ăn mừng thôi!''
Một ngày làm việc vất vả nhưng tràn ngập tiếng cười của gia đình Hikawa và các thành viên \hll\ đã kết thúc như thế đó.

\begin{center}
  {\ab BONUS GIF}
  \gif{2-8m}{1}{39}
\end{center}

\end{document}